% FILE: 02_lineage_and_context.tex
% Project: lifecycle
% Role: process-lifecycle-state-definitions-and-transitions

\section{Lineage and Context}
\label{sec:lifecycle-lineage}

\subsection{2016 Language-Model Lesson}
\label{sec:lifecycle-lineage-2016}

The foundational observation driving this research program is that local text imitation
does not conserve consequence. In 2016, a language model experiment demonstrated that a
system capable of producing locally plausible text---text that is fluent, contextually
appropriate, and indistinguishable from human-authored prose in a narrow window---could
nonetheless produce globally incoherent sequences: sequences whose later parts contradict
their earlier parts, whose claims do not follow from their premises, and whose action
recommendations conflict with the constraints established in the same document.

The lesson for process intelligence is direct. A process model that is locally
well-formed---each transition syntactically valid, each place correctly connected---may
be globally unsound: deadlocked at reachable markings, containing dead transitions that
consume tokens without enabling completion, or producing token accumulations that
violate boundedness. A governance report that is locally accurate---each metric correctly
computed from the data provided---may be globally unverifiable: the data provided may not
be the data that produced the operational outcome, and the metric may not have been
computed by the method stated.

The lifecycle framework operationalizes the 2016 lesson. It is not sufficient for a
process claim to be locally plausible. Every claim must resolve to a verifiable artifact
through a formally specified chain: model artifact, event log, conformance computation,
receipt. The typed evidence lifecycle $\mathcal{S}$ enforces this chain at the type level;
a value that has not passed through \texttt{Admit::admit()} cannot participate in
conformance computation, regardless of how plausible its content appears.

\subsection{Enterprise Process Gap}
\label{sec:lifecycle-lineage-gap}

The enterprise process gap identified in this research program is the structural
disconnection between what organizations claim about their processes and what their
event logs can prove. This gap is not a matter of bad faith; it is a structural property
of how enterprise process management has been practiced. Process models are maintained
as documentation artifacts, updated by human editors in response to audits and incidents.
Event logs, if collected at all, are maintained as compliance records in formats
(database exports, PDF audit reports, raw CSVs) that cannot be directly replayed against
the process model. The conformance check that would reveal the gap is not run, because
the infrastructure to run it does not exist in the standard enterprise stack.

The lifecycle framework names this gap precisely and provides the structural requirements
for closing it. A seller at process intelligence maturity L1 has no formal models and
no event log; their governance claims are narrative. A seller at L3 has formal models
and admitted OCEL or XES logs and can run periodic token replay; their governance
claims are partially verifiable. A seller at L5 has a full MAPE-K loop with receipted
artifacts at every stage and can replay any past loop cycle from the archived artifact
store; their governance claims are fully verifiable. The lifecycle project defines what
each maturity level requires, what each level can produce, and what audit responses
each level can honestly provide---drawn directly from \texttt{lifecycle/MATURITY\_MODEL.md}.

\subsection{Relationship to the Chatman Equation}
\label{sec:lifecycle-lineage-chatman}

The Chatman Equation, stated as $A = \mu(O^*)$, asserts that admissible process truth
is a function applied to the operational artifacts of the system under governance. The
lifecycle framework is the formal specification of what it means for an artifact to be
admissible ($A$), what constitutes the operational artifact set ($O^*$), and what the
function $\mu$ must compute to produce board-defensible governance claims.

In the lifecycle framework, an artifact is admissible if and only if it has traversed the
typed evidence lifecycle from \texttt{Raw} through \texttt{Admitted}. The operational
artifact set $O^*$ is the union of process models, event logs, conformance reports, repair
receipts, and decommissioning receipts accumulated across the full lifecycle. The function
$\mu$ is the composition of gate criteria: soundness verification at Gate 1, reachability
at Gate 2, fitness threshold at Gate 3, soundness preservation at Gate 4, discovery
conformance at Gate 5, and receipt archival at Gate 6. The equation $A = \mu(O^*)$ is
therefore not a metaphor; it is a computable function whose output is a set of
board-admissible statements backed by verifiable artifacts.

The Board Projection Holography framework, introduced in
\texttt{lifecycle/board-projection-holography.md}, extends the Chatman Equation to
prospective claims: a board projection is admissible if and only if it is derived from
the historical artifact set $O^*$ through the projection engine, with every projected
state carrying a cryptographic proof linking it to the underlying assumptions and the
historical baseline.

\subsection{Relationship to Receipts and Replay}
\label{sec:lifecycle-lineage-receipts}

The receipt and replay infrastructure is the execution layer of the lifecycle framework.
Every stage transition must produce a receipt. A repair action without a receipt is, by
definition of the lifecycle framework, not an execution for process intelligence
purposes---it is an unwitnessed intervention. A decommissioning without a cryptographic
decommissioning receipt $R_d$ is not an archival; it is a deletion. The receipt chain
is the mechanism by which the typed evidence lifecycle accumulates provenance across the
full operational history of a process model.

The BLAKE3 hash-chaining receipt ledger $H_n = \mathrm{BLAKE3}(\mathrm{Payload}_n
\parallel H_{n-1})$ ensures that the receipt chain is tamper-evident: any modification
to a past receipt invalidates all subsequent receipts in the chain. The Cryptographic
Decommissioning Receipt $R_d = \mathrm{Ed25519}_{K_{\mathrm{priv}}}(\mathrm{BLAKE3}(N)
\parallel \mathrm{BLAKE3}(L_{\mathrm{final}}) \parallel C_{\mathrm{total}} \parallel
F_{\mathrm{final}} \parallel T_{\mathrm{retire}})$ seals the lifecycle at decommission,
committing to the final model structure, the final event log, the total cost accumulated,
the final fitness score, and the retirement timestamp. The Oblivion Protocol---a
three-pass ChaCha20 CSPRNG wipe of WASM linear memory---ensures that the decommissioned
process model cannot be re-activated without re-traversing the full lifecycle from design.
